\documentclass[11pt]{article}

\usepackage[T1]{fontenc}
\usepackage{lmodern}
\usepackage{amsmath,amssymb,amsthm}
\usepackage[margin=1in]{geometry}
\usepackage{microtype}

\newtheorem{lemma}{Lemma}
\DeclareMathOperator{\arsinh}{arsinh}

\title{Analytic Positivity of the Aggregate Base Derivative}
\author{}
\date{}

\begin{document}
\maketitle

For \(0\leq s\leq1\), set
\[
 G_1(s)=\frac{\sqrt{1-s^2}-s\arccos s}{\pi}.
\]
For \(0<\rho<1\) and \(K>0\), put
\[
 \eta_\rho=G_1(\max\{\rho,1/2\}),\qquad
 d_\rho=\frac{2\arcsin\rho}{\pi},\qquad
 b_\rho=\frac{2\pi\rho}{1-\rho},
\]
\[
 \mu=\rho K,\qquad
 \overline R=\mu+\frac12,\qquad
 S=\sqrt{\overline R^2+b_\rho K},
\]
\[
 I(\rho,K)=\frac12\left[
 \overline R S+b_\rho K\arsinh
 \frac{\overline R}{\sqrt{b_\rho K}}-\overline R^2
 \right],
\]
and define the continuous aggregate reserve
\begin{equation}
\begin{aligned}
 \mathcal B(\rho,K)={}&
 (\mu-\tfrac12)(K\eta_\rho-1)
 +\frac{d_\rho}{2}(\mu-\tfrac32)(\mu-\tfrac12)\\
 &-(1+d_\rho)I(\rho,K)
 -\frac{8(\mu+\tfrac12)}{15\sqrt\mu}.
\end{aligned}
\label{cm:eq-B-def}
\end{equation}
Direct differentiation gives
\begin{equation}
 I_K=\rho(S-\overline R)
 +\frac{b_\rho}{2}\arsinh
 \frac{\overline R}{\sqrt{b_\rho K}},
\label{cm:eq-IK}
\end{equation}
and
\begin{equation}
\begin{aligned}
 \mathcal B_K={}&
 \rho(K\eta_\rho-1)+(\mu-\tfrac12)\eta_\rho
 +d_\rho\rho(\mu-1)\\
 &-(1+d_\rho)I_K
 -\frac{2\rho(2\mu-1)}{15\mu^{3/2}}.
\end{aligned}
\label{cm:eq-BK}
\end{equation}

% Analytic-replacement programme: standalone lemma.
% This file is deliberately independent of the interval certificates.
% Status: candidate proof, awaiting clean-room audit before insertion into the
% main manuscript.

\begin{lemma}[Analytic positivity of the aggregate base derivative]
\label{an:lem-BK200-positive}
For
\[
        \frac7{51}\leq \rho\leq\frac{39}{50},
\]
the continuous aggregate reserve defined in
\eqref{cm:eq-B-def} satisfies
\[
        \mathcal B_K(\rho,200)>0.
\]
This assertion uses no numerical or interval certificate.
\end{lemma}

\begin{proof}
We use only the elementary rational bounds
\begin{equation}
  3<\pi<\frac{22}{7},\qquad
  \frac{78}{25}<\pi,
  \qquad
  \frac75<\sqrt2,
  \qquad
  \frac{265}{153}<\sqrt3<\frac{26}{15}.
  \label{an:eq-elementary-constants}
\end{equation}
The two square-root bounds in \eqref{an:eq-elementary-constants} follow
immediately by squaring; the displayed rational bounds for \(\pi\) are the
classical Archimedean bounds.

Put
\[
  \mu=200\rho,
  \qquad
  b_\rho=\frac{2\pi\rho}{1-\rho},
  \qquad
  \overline R=200\rho+\frac12,
  \qquad
  S=\sqrt{\overline R^2+200b_\rho},
\]
and
\[
  u=\frac{\overline R}{\sqrt{200b_\rho}},
  \qquad P=S-\overline R.
\]
We first obtain a uniform elementary upper bound for the derivative of the
integral term.  Since \(\rho>1/8\), one has
\(\overline R<204\rho\).  Therefore
\[
 u^2
 <\frac{204^2\rho(1-\rho)}{400\pi}
 \leq\frac{41616}{1600\pi}
 <\frac{867}{100}<9.
\]
Thus \(u<3\).  Moreover
\[
 \sinh 2=2+\frac{2^3}{3!}+\frac{2^5}{5!}+\cdots
 >\frac{10}{3}>3,
\]
and hence \(\operatorname{arsinh}u<2\).  The identity
\[
 P=\frac{200b_\rho}{S+\overline R}
\]
and the inequalities \(S>\overline R>200\rho\) give
\(P<b_\rho/(2\rho)\).  Formula \eqref{cm:eq-IK} now yields
\begin{equation}
 I_K(200)
 =\rho P+\frac{b_\rho}{2}\operatorname{arsinh}u
 <\frac{3b_\rho}{2}.
 \label{an:eq-IK-simple-upper}
\end{equation}

The final error term in \eqref{cm:eq-BK} obeys
\begin{equation}
 0<\frac{2\rho(2\mu-1)}{15\mu^{3/2}}
 <\frac{4\sqrt\rho}{15\sqrt{200}}
 <\frac{2}{105}<\frac1{50}.
 \label{an:eq-BK-error-upper}
\end{equation}
At \(K=200\), formula \eqref{cm:eq-BK} becomes
\begin{equation}
\begin{split}
 \mathcal B_K(\rho,200)
 ={}&400\rho\eta_\rho-\rho-\frac{\eta_\rho}{2}
 +d_\rho\rho(200\rho-1)\\
 &-(1+d_\rho)I_K(200)
 -\frac{2\rho(400\rho-1)}{15(200\rho)^{3/2}}.
\end{split}
\label{an:eq-BK200-expanded}
\end{equation}

We first treat \(7/51\leq\rho\leq1/2\).  On this interval
\[
 \eta_\rho=\omega_0=\frac{\sqrt3}{2\pi}-\frac16,
 \qquad 0<d_\rho\leq\frac13.
\]
The last two bounds in \eqref{an:eq-elementary-constants} give
\[
 \frac{13}{120}<\eta_\rho<\frac19.
\]
Indeed,
\[
 \frac{\sqrt3}{2\pi}>
 \frac{265}{153}\frac7{44}=\frac{1855}{6732}>
 \frac{11}{40},
\]
whereas
\[
 \frac{\sqrt3}{2\pi}<
 \frac{26}{15}\frac{25}{156}=\frac5{18}.
\]
The term \(d_\rho\rho(200\rho-1)\) is nonnegative.  Combining
\eqref{an:eq-IK-simple-upper}, \eqref{an:eq-BK-error-upper},
\(2\pi<44/7\), and \eqref{an:eq-BK200-expanded}, we obtain
\begin{align*}
 \mathcal B_K(\rho,200)
 &>\rho\left(\frac{127}{3}
       -\frac{88}{7(1-\rho)}\right)-\frac{17}{225}\\
 &\geq \frac7{51}\left(\frac{127}{3}-\frac{176}{7}\right)
       -\frac{17}{225}
 =\frac{2912}{1275}>0.
\end{align*}

It remains to treat \(1/2\leq\rho\leq39/50\).  Monotonicity of
\(\arcsin\), together with \(\arcsin\rho\geq\rho\), gives
\begin{equation}
       \frac{7\rho}{11}<d_\rho<\frac35.
 \label{an:eq-d-high-bounds}
\end{equation}
For the upper bound, note that
\[
 \frac{39}{50}<\sin\frac{3\pi}{10}
 =\frac{1+\sqrt5}{4};
\]
the strict rational comparison follows from
\(\sqrt5>53/25\).  For the lower bound use
\(2/\pi>7/11\).

Since
\[
 \eta_\rho=G_1(\rho)
 =\frac1\pi\int_\rho^1\arccos t\,dt
\]
and \(\arccos t\geq\sqrt{2(1-t)}\), we have
\begin{equation}
 \eta_\rho
 \geq\frac{2\sqrt2}{3\pi}(1-\rho)^{3/2}
 >\frac{49}{165}(1-\rho)^{3/2}.
 \label{an:eq-eta-high-lower}
\end{equation}
Also \(0<\eta_\rho\leq\omega_0<1/9\).

Let \([a,c]\) be any of the six intervals in the table below, and let
\(s_c\) be the displayed rational number.  Direct squaring proves
\(s_c^2\leq1-c\).  Hence \eqref{an:eq-eta-high-lower} gives
\[
 \eta_\rho\geq \eta_c
 :=\frac{49}{165}(1-c)s_c
 \qquad(a\leq\rho\leq c).
\]
Equations \eqref{an:eq-IK-simple-upper}--\eqref{an:eq-d-high-bounds}
then imply the entirely rational lower bound
\begin{equation}
 \mathcal B_K(\rho,200)>L(a,c),
 \label{an:eq-BK-slab-bound}
\end{equation}
where
\[
 L(a,c):=
 400a\eta_c-c-\frac1{18}
 +\frac{7a^2}{11}(200a-1)
 -\frac{528}{35}\frac{c}{1-c}-\frac1{50}.
\]
Every entry in the last column is obtained by substitution and integer
cross-multiplication:
\[
\begin{array}{c|c|c|c|c}
 a&c&s_c&\eta_c&L(a,c)\\ \hline
 \frac12&\frac{11}{20}&\frac23&\frac{49}{550}
   &\frac{251291}{17325}\\[2pt]
 \frac{11}{20}&\frac35&\frac58&\frac{49}{660}
   &\frac{70619}{5040}\\[2pt]
 \frac35&\frac{13}{20}&\frac7{12}&\frac{2401}{39600}
   &\frac{576451}{44100}\\[2pt]
 \frac{13}{20}&\frac7{10}&\frac6{11}&\frac{147}{3025}
   &\frac{4940821}{435600}\\[2pt]
 \frac7{10}&\frac34&\frac12&\frac{49}{1320}
   &\frac{2411}{315}\\[2pt]
 \frac34&\frac{39}{50}&\frac{23}{50}&\frac{1127}{37500}
   &\frac{11101801}{1386000}
\end{array}
\]
All six rational lower bounds are strictly positive, and the intervals
tile \([1/2,39/50]\).  This proves the assertion on the high branch and
completes the proof.
\end{proof}


\end{document}
